% from ega3-6-fr.tex
\oldpage[III]{166}
\subsection{Les suites spectrales de changement de base dans les hypertor
globaux.}

\begin{env}[6.9.1]
Les hypothèses et notations étant toujours celles de (6.7.11) (et en
particulier les $\mathscr P_\bullet^{(i)}$ étant supposés limités
inférieurement), considérons un morphisme $g:S'\to S$ de préschémas, et
posons
$Y'^{(i)}=Y_{(S')}^{(i)}$, $X'^{(i)}=X_{(S')}^{(i)}$ et
$\mathscr P_\bullet'^{(i)}=
\mathscr P_\bullet^{(i)}\otimes_{\mathscr O_S}\mathscr O_{S'}$,
$\mathscr P_\bullet'^{(i)}$ étant donc un complexe de
$\mathscr O_{X'^{(i)}}$-Modules quasi-cohérents; soit
$f'_i=(f_i)_{(S')}:X'^{(i)}\to Y'^{(i)}$, qui est un morphisme séparé et
quasi-compact (I, 5.5.1 et 6.6.4). Nous nous proposons d'étudier les relations
entre les $\mathscr O_{Y'}$-Modules quasi-cohérents
$\mathfrak{Tor}_n^{S'}((f'_i)_{i\in I};
(\mathscr P_\bullet'^{(i)})_{i\in I})$ et
$\mathfrak{Tor}_n^S((f_i)_{i\in I};
(\mathscr P_\bullet^{(i)})_{i\in I})
\otimes_{\mathscr O_S}\mathscr O_{S'}$, où
$Y'=Y\times_S S'=Y_{(S')}$. Un cas particulièrement simple est le suivant,
qui se réduit à (1.4.15) lorsque $I$ se réduit à un seul élément et
$\mathscr P_\bullet$ à un seul module :
\end{env}

\begin{proposition}[6.9.2]
Si le morphisme $g:S'\to S$ est plat, on a un isomorphisme canonique de
$\partial$-foncteurs (en les $\mathscr P_\bullet^{(i)}$) :
\[
  \mathfrak{Tor}_\bullet^S
  ((f_i)_{i\in I};(\mathscr P_\bullet^{(i)})_{i\in I})
  \otimes_{\mathscr O_S}\mathscr O_{S'}
  \xrightarrow{\sim}
  \mathfrak{Tor}_\bullet^{S'}
  ((f'_i)_{i\in I};(\mathscr P_\bullet'^{(i)})_{i\in I}).
  \tag{6.9.2.1}\label{III.6.9.2.1.fr}
\]
\end{proposition}

\begin{proof}
On peut encore se borner au cas où $S$, $S'$ et les $Y^{(i)}$ sont affines,
le recollement se faisant suivant les méthodes de (6.7.1) et (6.7.2). Soient
$S=\operatorname{Spec}(A)$, $S'=\operatorname{Spec}(A')$, et prenons pour
chaque $i$ un recouvrement ouvert affine $\mathfrak U^{(i)}$ de $X^{(i)}$;
si $u_i:X'^{(i)}\to X^{(i)}$ est la projection canonique,
$u_i^{-1}(\mathfrak U^{(i)})$ est un recouvrement ouvert affine de
$X'^{(i)}$ (II, 1.5.5), que nous noterons $\mathfrak U'^{(i)}$; il est clair
alors que
\[
  C^\bullet(\mathfrak U'^{(i)},\mathscr P_\bullet'^{(i)})
  =C^\bullet(\mathfrak U^{(i)},\mathscr P_\bullet^{(i)})\otimes_A A',
\]
et l'existence de l'isomorphisme (6.9.2.1) est immédiate, car si
$M_{\bullet\bullet\bullet}^{(i)}$ est une résolution projective de
Cartan-Eilenberg de
$L_{\bullet\bullet}^{(i)}=C^\bullet(\mathfrak U^{(i)},
\mathscr P_\bullet^{(i)})$ au sens de (0, 11.7.1), formée de $A$-modules,
$M_{\bullet\bullet\bullet}^{(i)}\otimes_A A'$ est une résolution projective
de Cartan-Eilenberg (au même sens) de
$L_{\bullet\bullet}^{(i)}\otimes_A A'$ formée de $A'$-modules, en vertu de
l'hypothèse que $A'$ est un $A$-module plat; cette même hypothèse montre en
outre que
\[
  H_\bullet\left(
    \bigotimes_{i=1}^m
    (M_{\bullet\bullet\bullet}^{(i)}\otimes_A A')
  \right)
  =H_\bullet\left(
    \bigotimes_{i=1}^m M_{\bullet\bullet\bullet}^{(i)}
  \right)\otimes_A A'.
\]
\end{proof}

\begin{env}[6.9.2.2]
On notera que lorsque $I$ est réduit au seul élément $1$, la formule
(6.9.2.1) se déduit directement de (6.7.7), appliqué en prenant
$Y_2=X_2=S'$, $f_2=1_{S'}$, et le complexe $\mathscr P_\bullet^{(2)}$
réduit à son terme de degré $0$, égal à $\mathscr O_{S'}$; on sait alors que
l'hypercohomologie $\mathscr H^n(1_{S'},\mathscr O_{S'})$ est nulle pour tout
$n\ne0$ et se réduit à $\mathscr O_{S'}$ pour $n=0$ (6.2.1).
\end{env}

Dans le cas général, nous allons introduire à la place de $\mathscr O_{S'}$
un complexe $\mathscr Q'_\bullet$ de $\mathscr O_{S'}$-Modules
quasi-cohérents limités inférieurement, de sorte que si, pour simplifier, on
prend $I=\{1,2,\ldots,m\}$, on peut considérer le $\partial$-foncteur
\[
  \mathfrak{Tor}_\bullet^S
  (f_1,\ldots,f_m,1_{S'};
   \mathscr P_\bullet^{(1)},\ldots,\mathscr P_\bullet^{(m)},
   \mathscr Q'_\bullet).
\]

\begin{proposition}[6.9.3]
Il existe trois foncteurs spectraux canoniques biréguliers notés
${}^{(t)}\mathscr E(f_1,\ldots,f_m;
\mathscr P_\bullet^{(1)},\ldots,\mathscr P_\bullet^{(m)},
\mathscr Q'_\bullet)$ (avec $t=e$, $f$ ou $f'$) ayant pour aboutissement
commun
$\mathfrak{Tor}_\bullet^S(f_1,\ldots,f_m,1_{S'};
\mathscr P_\bullet^{(1)},\ldots,\mathscr P_\bullet^{(m)},
\mathscr Q'_\bullet)$ et dont les termes $E_2$ sont respectivement

